\documentclass{internal-report}
\usepackage[UTF8]{ctex}
\usepackage{booktabs}
\usepackage{multirow}

% STATUS: draft
% LAST_MODIFIED: 2026-08-07

\title{C 题初始分析：面向算电协同的多目标调度优化研究}
\author{MathModel 建模小组}
\date{2026-08-07}

\begin{document}

\maketitle

\section{问题重述}

在"算电协同"国家战略背景下（2026 年政府工作报告首次纳入国家级新基建工程），
本研究针对六个典型区域（东部高负荷 RegionA/B/C、西部算力中心 RegionD、
光伏富集 RegionE、风电富集 RegionF）中的 AI 任务调度与电力系统协同优化问题。

系统主运行时域为第 0--2399 小时，第 2400--2405 小时结清末端任务，
第 2406 小时仅做电力与储能状态结算。任务分三类：实时推理（到达即开工、时延敏感）、
批量推理（中弹性）、AI 训练（低时延敏感，为低碳调度的主要调节对象）。
任务不可抢占、不可拆分、不可中途迁移，均须在时点 2406 前完成。

\section{子问题拆解}

\begin{table}[h]
\centering
\footnotesize
\begin{tabular}{p{0.7cm}p{3.4cm}p{3.4cm}p{4.4cm}}
\toprule
编号 & 子问题 & 决策变量 & 目标/指标 \\
\midrule
S1 & GPU 需求统计分析与短期预测 + 基础算力调度 &
  任务区域分配、开工时段 & 预测精度、GPU 利用率、按时完成 \\
S2 & 碳感知任务调度模型 &
  任务迁移目的地、开工时段 & 运行成本、碳排放、网络时延、新能源利用率 \\
S3 & 储能协同优化模型 &
  储能充放电、购售电、新能源分配 & 成本、碳排、峰值净购电、负荷波动 \\
S4 & 多区域算--储--电协同优化模型 &
  S2+S3 全部变量 & 成本、碳排、时延、QoS、新能源利用率、峰值净购电 \\
\bottomrule
\end{tabular}
\caption{子问题拆解}
\label{tab:subproblems}
\end{table}

四个子问题逐步放开自由度：S1 为基础调度（预测驱动），S2 引入任务迁移与开工时段自由度，
S3 固定任务调度只优化储能侧，S4 统一优化算--储--电全系统并进行多目标权衡与场景比较。

\section{相关知识}

\subsection{优化类知识卡}

\begin{itemize}
  \item MS-018（线性/混合整数规划）：大规模排程精确优化，适用于 S1/S2 约束可线性化的情形。
  \item MS-016（NSGA-II）：混合多目标优化，适用于 S4 六目标 Pareto 权衡。
  \item MS-030（双层随机规划）：层级决策分离目标优先级，可类比算力中心与储能运营的分层。
  \item AL-003（遗传算法）、PR-001（约束处理优先级：编码 $>$ 罚函数 $>$ 修复）、
        WF-001（优化求解工作流：建模 $\rightarrow$ 启发式快速解 $\rightarrow$ 精确校验）。
\end{itemize}

\subsection{预测类知识卡}

\begin{itemize}
  \item MS-009（Prophet + 外生变量）、MS-015（LSTM 多周期时间序列）、AL-006（ACF 周期分析）。
  \item PF-003：ARIMA 在非平稳不规则序列上的失败教训——预测前必须先做平稳性与周期性诊断。
\end{itemize}

\subsection{方案对比经验}

SP-003（2024C 优化哲学）：精确 LP 与启发式（GA/DEGA）的选型取决于约束是否可线性化、
是否需要精确最优保证、是否需要灵活编码复杂结构。本题 S2/S4 的大规模离散调度问题
（约 5 万任务 × 2400 小时）需权衡求解规模，倾向时间分解 + 精确求解器或启发式路线。

\subsection{领域知识（来自阶段 0.0）}

算电协同的核心机制是负载时移与地理迁移：东部高电价高碳（A 电价 468 元/MWh、碳强度
0.668 tCO$_2$/MWh）vs 西部低电价低碳（E 电价 247 元/MWh、碳强度 0.237 tCO$_2$/MWh）；
RegionD 算力容量最大（1600 GPU）。储能 SOC 递推式与统一口径公式已在
\texttt{solution/domain-knowledge.md} 中固化。

\section{初始建模假设}

\begin{enumerate}
  \item 任务到达序列在 0--2399 小时全部已知（确定性输入）；问题二起按实际到达计算。
  \item S1"基础调度"暂假定不引入跨区域迁移（或仅训练任务可跨区），待方案探索确认。
  \item 电价、碳强度、可用新能源为逐时已知输入；不建模预测误差（S4 场景分析另行处理）。
  \item 问题一、二不考虑储能；问题三、四考虑储能（问题三固定任务调度）。
  \item QoS 服务质量需在方案探索阶段给出量化定义（如按时完成率、时延满意度）。
  \item 迁移决策仅受时延 SLA 与容量/电力约束影响（不建模带宽、传输能耗与潮流）。
\end{enumerate}

\section{待解决的开放问题}

\begin{enumerate}
  \item S1 中"基础调度模型"是否允许跨区域迁移？实时任务是否强制本地执行？
  \item S2 调度模型求解规模：5 万任务 × 2400 小时的 MILP 需时间分解还是启发式？
        多目标四指标如何权衡（加权和 vs $\varepsilon$-约束 vs NSGA-II）？
  \item S3 储能优化目标函数的具体形式（成本+碳排加权系数如何标定）？
        负荷波动程度的量化指标（峰谷差 / 方差 / 均方差）？
  \item S4 六目标权衡方法选择与 QoS 度量定义；场景比较的维度与基准场景设定。
  \item 新能源利用率与基准（附件 SOC\_MWh 轨迹）的对比口径。
\end{enumerate}

\end{document}
